\section{Limitations and future work}

The uniform squared-risk comparison in \cref{obj:thm:observed-record-precision-frontier} characterizes attainable precision for a single experiment under the schedule restrictions in \cref{obj:def:model} and the observation design in \cref{obj:ass:assignment,obj:ass:audit,obj:ass:design-independence}. Its comparison constants establish the minimax order across the admissible population sizes, degree bounds, and retention probabilities. Exact minimax constants remain a separate question, as does uncertainty quantification with specified confidence-interval coverage. A squared-risk guarantee controls expected estimation error; constructing informative intervals requires an additional coverage argument that accounts for assignment variation and sampled graph information. The behavior of such intervals near regimes in which the risk remains bounded away from zero warrants particular attention.

Known retention is central to the inverse-inclusion score in \cref{obj:def:audit-score}. When retention probabilities are estimated from calibration data, their estimation error enters the weights and can affect both bias and variance, especially at low retention. An extension would need to specify the calibration experiment, its relationship to the treatment experiment, and conditions ensuring sufficiently accurate estimation of inclusion probabilities. Heterogeneous retention probabilities present a related problem: calibration and precision guarantees would have to account for which arrows are sampled less frequently and how their spillover coefficients contribute to the target.

The audit design also distinguishes missing arrows from falsely recorded arrows. Under \cref{obj:ass:audit}, the recorded graph is formed by retaining true arrows independently. Allowing false-positive edges changes the observation law and the behavior of the score. A useful extension would specify separate inclusion and false-positive mechanisms, identify which aspects of those mechanisms are known or estimable, and determine how their uncertainty affects total-effect estimation.

Connecting the calibrated audit experiment to a platform collection mechanism requires calibration at the level of directed interference arrows. For the present guarantees to apply, the collection procedure must induce the common known inclusion probability, independent arrow retention, and independence from treatment assignment required by \cref{obj:ass:audit,obj:ass:design-independence}. It must also record assignments and realized outcomes for the full study population, while the interference schedule satisfies \cref{obj:def:model}. Sampling interaction events can induce dependent or heterogeneous arrow inclusion when several events represent the same arrow. Translating such a workflow into the present experiment therefore requires an explicit mapping from event records to interference arrows and justification of the resulting inclusion law. Treatment-dependent activity or outcome-dependent recording requires a different observation model.

Alternative assignments and repeated experiments offer further directions. Unequal treatment probabilities, complete randomization, and cluster assignment alter both the score construction and the dependence relevant to precision bounds. Repeated experiments would require specifying whether the graph and response coefficients remain fixed, how assignments vary across rounds, and whether graph audits are shared or renewed. Their potential precision gains depend on these choices. Richer response models, including treatment interactions and nonlinear exposure responses, likewise require an expanded schedule class and restrictions controlling how joint treatment changes affect outcomes. The additive coefficient-mass restriction in \cref{obj:ass:coefficient-mass} provides the starting point for such extensions.

Graph recovery is a separate inferential objective from estimating the population total treatment effect. Recovering individual arrows requires a recovery criterion and identifying restrictions connecting arrow presence to observable responses, together with an observation design that supplies sufficient information about individual associations. Studying those requirements alongside total-effect precision would clarify when collection sufficient for a population contrast also supports learning the interference structure.
